\section{Introduction}
Sharing information is vital for communication and discourse across domains, as it allows for knowledge to be disseminated to a wider audience. This is often done by users through documents in multiple formats that nevertheless share some underlying knowledge. A product manager may need to create a requirements spec, a product pitch deck, and an announcement newsletter for the same project. Likewise, a person on the job market may create a resume, a cover letter, and a personal website. We consider these documents to be \textit{templatic views} of the same underlying knowledge.

This is equally true for the scientific domain, in which researchers create documents in multiple formats to effectively communicate and showcase their work, -- such as through academic papers, conference talks, social media posts, poster presentations, and non-technical blog posts. Sharing knowledge in multiple formats broadens the audience and can help bridge the information gap between domain experts, researchers in adjacent fields, and even the general public, leading to greater understanding, collaborations and accelerated progress~\cite{Bornmann2014GrowthRO}.


Past work on document generation has focused on developing generation and evaluation methods specific to a single document type~\cite{Fu2021DOC2PPTAP,Qiang2016LearningTG,chandrasekaran-etal-2020-overview-insights}. 
Narrow, custom methods tailored to individual document types are, nevertheless, time consuming to engineer and manage over the long term. For example, in an enterprise setting, it’s common to have dozens of occupation- and task-specific documents, each with their own template.
Additionally, specific trained methods require data that may be expensive to acquire, or even be unavailable entirely. Meanwhile, LLMs have recently shown great success in long document generation~\cite{Radford2019LanguageMA, Brown2020LanguageMA}, indicating that this fragmentation of methods may no longer be necessary. Thus, our goal is to unify methods for both generating and evaluating templatic views of documents, allowing system designers and engineers to manage and adapt to a range of document types and domains easily and efficiently. 

We begin by showing that LLMs are capable of diverse, structured document generation, requiring very little instructional guidance to do so effectively. Additionally, a few minor augmentations to the prompt -- such as a structured, intermediate representation, and simple stylistic descriptions -- can further improve downstream performance, especially for smaller, less resource intensive models. These findings have important implications on the deployment and scaling of unified, real-world AI-assisted document authoring systems.

In similar vein, we then introduce Template Adaptable Evaluation (TAE), departing from prior work's task specific evaluation methods~\cite{Zhang2019BERTScoreET, Qiang2016LearningTG,Wang2015iPosterIP}. TAE is a unified precision-recall style framework for automatic evaluation that is highly customizable, allowing users to easily integrate existing text-based metrics from the literature into its formulation and tailor it to their specific use case.% This adaptability includes styles of documents for which there is no publicly available data, such as enterprise-internal product documents. 
Additionally, this framework allows developers to compare performance across document types, without needing to develop an evaluation metric for each individual template.


We evaluate our unified approach for templatic view generation and evaluation on 3 types of documents: slides, posters, and blog posts~\cite{Fu2021DOC2PPTAP,Qiang2016LearningTG,chandrasekaran-etal-2020-overview-insights}. %We use the DOC2PPT~\cite{Fu2021DOC2PPTAP}, Paper-Poster~\cite{Qiang2016LearningTG}, and LongSumm~\cite{chandrasekaran-etal-2020-overview-insights} datasets, respectively, where the inputs are scientific articles and the outputs are the documents of the corresponding types.
Our experiments demonstrate that using a structured intermediate representation leads to improvements in performance across tasks, with greater gains for smaller language language models.
%demonstrating that LLMs are capable of generating templatic views with structure-augmented prompts, despite having no supervision.
In our human evaluation to validate both our unified document generation method and evaluation metric, we show that annotators prefer the output yielded by the structure-aware generation process 82\% of the time and that our evaluation metric correlates more highly with human preference than other popular metrics. A visualization of our methods can be found in Figure~\ref{fig:method}. We release our code\footnote{\url{https://github.com/microsoft/knowledge-centric-templatic-views}} to support future research.
